% ---------------------------------------------------------------------------
% Resumen (norma 3.4)
% ---------------------------------------------------------------------------
% Requisitos de la norma:
%   - Máximo UNA página.
%   - Típicamente un SOLO PÁRRAFO de 250 palabras como máximo.
%   - Debe exponer: tema tratado, metodología empleada, resultados obtenidos
%     y conclusiones alcanzadas.
%   - Si se usan siglas, su significado se explica entre paréntesis dentro del
%     propio resumen, la primera vez que aparezcan.
% Redactado al cierre de la Fase 6, cuando ya existen los resultados medidos.
% 250 palabras exactas.

\seccionSinNumerar{Resumen}

Los sistemas de seguridad basados en reglas generan más alertas de las que un equipo puede
analizar, muchas falsas y ninguna justificada. Este trabajo, realizado como pasantía larga en
Domotai C.A., desarrolla un motor de triaje que clasifica y prioriza alertas, justifica cada
decisión con un modelo de lenguaje (LLM) apoyado en generación aumentada por recuperación (RAG)
en local, y propone una acción de respuesta contenida sujeta a validación
humana: el modelo justifica, no decide. Se modeló un cliente genérico, se revisó el estado del arte, se construyó un
laboratorio reproducible con Containerlab y Wazuh, se diseñó una arquitectura que separa
clasificación, política determinista y perfil de cliente configurable, y se evaluó el prototipo
frente al nivel de regla de Wazuh. Sobre 300 alertas reales de tres familias,
incluidos casos diseñados para que las reglas fallaran, el prototipo redujo la tasa de falsos
positivos a 0,009 frente a 0,310 del método por reglas y detectó todas las amenazas, incluidas las
que el nivel de regla deja pasar (exhaustividad 0,920 del baseline), sin acciones disruptivas. Un árbol de decisión entrenado sobre esos datos redescubrió las reglas escritas a mano
y aportó una, la de la ráfaga, adoptada como regla auditable. La recuperación aumentada corrigió
la imprecisión de las explicaciones. Queda como limitación el clasificador con ajuste fino, que
exige más variedad de datos de la que el laboratorio produce.
\vspace{1.5\baselineskip}

\noindent\textbf{Palabras clave:} triaje de alertas de seguridad, modelos de lenguaje de gran
escala, generación aumentada por recuperación, MITRE ATT\&CK, validación humana.
